\section{Discussão e Conclusão}

O dimensionamento preliminar indica que é possível obter uma configuração equivalente ao Exocet MM40 Block 2, com \(70~\text{km}\) de alcance e Mach \(0{,}9\), usando apenas propelente sólido, desde que seja adotado um motor de queima prolongada com empuxo médio de aproximadamente \(2{,}5~\text{kN}\).

A principal consequência da substituição de uma arquitetura com sustentação eficiente por uma solução exclusivamente sólida é a necessidade de uma fração de propelente relativamente alta, estimada em \(28\%\) da massa total. Isso reduz a margem disponível para estrutura, sistemas, carga útil e integração.

O resultado final adotado possui \(850~\text{kg}\), \(5{,}8~\text{m}\) de comprimento, \(0{,}35~\text{m}\) de diâmetro e \(240~\text{kg}\) de propelente sólido. A configuração fecha o impulso total requerido com margem de \(20\%\), mas deve ser entendida como um dimensionamento preliminar. Um projeto detalhado exigiria refinamento aerodinâmico, cálculo de trajetória, distribuição de massa, estabilidade, controle, análise térmica e dimensionamento propulsivo específico.

\section*{Referências}

\begin{itemize}
    \item CSIS Missile Threat. \textit{Exocet}. Disponível em: \url{https://missilethreat.csis.org/missile/exocet/}.
    \item GlobalSecurity. \textit{Exocet AM.39 / MM.40 Specifications}. Disponível em: \url{https://www.globalsecurity.org/military/world/europe/exocet-specs.htm}.
    \item WeaponSystems.net. \textit{MM40 Exocet}. Disponível em: \url{https://weaponsystems.net/system/1126-MM40+Exocet}.
\end{itemize}
